\documentclass{beamer}
\usepackage{microtype}
\usepackage{default}

\usetheme{simple}
\usepackage{fontspec}
\usepackage{graphicx}
\usepackage[utf8]{inputenc}
\usepackage[justification=centering]{caption}
\usepackage{subcaption}
\setmainfont{Fira Sans}
\setbeamertemplate{caption}{\raggedright\insertcaption\par}
\captionsetup[subfigure]{labelformat=empty}
\newfontfamily\DejaSans{DejaVu Sans}
\setbeamerfont{title}{family=\texttt,size=\huge}
\usepackage[scale=2]{ccicons}
\title{The i-score interactive sequencer}
\subtitle{an intermedia sequencer for interactive scenarios authoring}
\date{\today}
\author{Jean-Michaël Celerier, Théo de la Hogue}
\institute{LaBRI, Blue Yeti, GMEA }
\date{January 30, 2016}

\begin{document}
    
\maketitle

\begin{frame}
    \frametitle{The problem}        
\end{frame}

\begin{frame}
    \frametitle{Example installations}        

    \begin{figure}
        \centering
        \includegraphics[width=0.9\textwidth]{images/futuroscope.jpg}
        \caption[]{Futuroscope, France. Blue Yeti}
    \end{figure}
\end{frame}


\begin{frame}
    \frametitle{Example installations}        
    robots
    ensatt
    pascal
    
\end{frame}

\begin{frame}
    \frametitle{Screenshot}        
\end{frame}


\begin{frame}
    \frametitle{Contributors, Companies, Agencies involved}
    \begin{figure}[htbp]
    \begin{subfigure}[b]{0.25\textwidth} 
        \centering
        \begin{tikzpicture}[scale=.05]

\pgfmathsetmacro\bleulabriR{0/255}
\pgfmathsetmacro\bleulabriG{114/255}
\pgfmathsetmacro\bleulabriB{185/255}

\pgfmathsetmacro\vertlabriR{10/255}
\pgfmathsetmacro\vertlabriG{177/255}
\pgfmathsetmacro\vertlabriB{134/255}

\pgfmathsetmacro\rougelabritopR{162/255}
\pgfmathsetmacro\rougelabritopG{47/255}
\pgfmathsetmacro\rougelabritopB{50/255}

\pgfmathsetmacro\rougelabrimiddleR{178/255}
\pgfmathsetmacro\rougelabrimiddleG{117/255}
\pgfmathsetmacro\rougelabrimiddleB{99/255}

\pgfmathsetmacro\rougelabribottomR{230/255}
\pgfmathsetmacro\rougelabribottomG{206/255}
\pgfmathsetmacro\rougelabribottomB{196/255}


\definecolor{bleulabri}{rgb}{\bleulabriR, \bleulabriG, \bleulabriB}
\definecolor{vertlabri}{rgb}{\vertlabriR, \vertlabriG, \vertlabriB}
\definecolor{rougetoplabri}{rgb}{\rougelabritopR, \rougelabritopG,
 \rougelabritopB}
\definecolor{rougemiddlelabri}{rgb}{\rougelabrimiddleR,
 \rougelabrimiddleG, \rougelabrimiddleB}
\definecolor{rougebottomlabri}{rgb}{\rougelabribottomR,
 \rougelabribottomG, \rougelabribottomB}

\def\rectanglepath{-- ++(0cm,38.5cm)-- ++(10cm,0cm)--
 ++(0cm,-30.3cm)-- ++(13cm,0cm)-- ++(-2cm,-8.2cm)-- cycle}


\shade[top color=rougetoplabri!95!black,bottom
 color=rougebottomlabri!15!white!95!black,middle
 color=rougemiddlelabri!90!white!98!black,shading=axis,
shading angle=-135] (19.25, 19.25) circle (6.75cm);

\fill[fill=bleulabri] (0,0) \rectanglepath;
\fill[rotate=180,fill=vertlabri] (-38.5,-38.5) \rectanglepath;

\end{tikzpicture}

        \caption{LaBRI \\ \url{www.labri.fr}}
    \end{subfigure}
    ~
    \begin{subfigure}[b]{0.25\textwidth} 
        \centering
        \includegraphics[scale=0.12]{images/by.png}
        \caption[]{Blue Yeti\\ \url{www.blueyeti.fr}}
    \end{subfigure}
    ~
    \begin{subfigure}[b]{0.25\textwidth} 
        \centering
        \includegraphics[scale=0.15]{images/gmea.png}
        \caption[]{GMEA\\\url{www.gmea.net}}
    \end{subfigure}
    
    
    \begin{subfigure}[b]{0.25\textwidth} 
        \centering
        \includegraphics[scale=0.25]{images/cnam.png}
        \caption[]{CNAM : \\ CEDRIC, ENJMIN\\ \url{cedric.cnam.fr}}
    \end{subfigure}
    ~
    \begin{subfigure}[b]{0.25\textwidth} 
        \centering
        \includegraphics[scale=0.05]{images/ists.jpg}
        \caption[]{ISTS\\ \url{cedric.cnam.fr}}
    \end{subfigure}
    ~
    \begin{subfigure}[b]{0.25\textwidth} 
        \centering
        \includegraphics[scale=0.23]{images/ensatt.jpg}
        \caption[]{ENSATT\\ \url{cedric.cnam.fr}}
    \end{subfigure}    
\end{figure} 
    
    {\Large Artists: Les Baltazars, Renaud Rubiano, Antoine Villeret...}
    
\end{frame}

    \begin{frame}
        \frametitle{What i-score is : }
        \begin{itemize}        
        \item Visual programming language for interactive multimedia applications        
        \item Free software : GPL v3 (UI) \& LGPL v2.1 (Engine)
        \item C++ (Qt, CMake)        
        \item Linux / OS X / Windows        
        \item Alpha-quality \DejaSans{☹}
        \end{itemize}
    \end{frame}
    \begin{frame}
        \frametitle{What i-score is not : }
        \begin{itemize}
            \item Audio sequencer        
            \item Video sequencer        
            \item PureData or Ableton Live       
            \item Bug-free
        \end{itemize}
    \end{frame}
    
    \begin{frame}
        \frametitle{Inter-operability}
        Max, PureData, Unity, OpenFrameworks, Processing, Jamoma, Modul8, Millumin, Quartz Composer, Qt...
    \end{frame}
    
    \begin{frame}
        \centering \Huge Demonstration
    \end{frame}
    
    \begin{frame}
        \frametitle{Automations, mappings}
        
    \end{frame}
    
    \begin{frame}
        \frametitle{Javascript}
        
    \end{frame}
    
    \begin{frame}
        \frametitle{Hierarchy}
        
    \end{frame}
    
    \begin{frame}
        \frametitle{Spatial automations}
        
    \end{frame}
    
    \begin{frame}
        \frametitle{Future : audio ?}
        
    \end{frame}

\begin{frame}
    \frametitle{Future : distribution ?}
    
\end{frame}

\begin{frame}
    \frametitle{Future : other features}
    
    \begin{itemize}
    \item MIDI, WebSockets support
    \item Some level of patching, like Pd
    \item Complete remote-control abilities; currently execution can be followed via a web page.
    \end{itemize}
\end{frame}

\begin{frame}
    \frametitle{Contributing}    
    \begin{itemize}
    \item UX, UI (mock-ups were done but not entirely implemented)
    \item Documentation, writing demo scenarios
    \item Translations    
    \item Implement the Minuit protocol in your software with the OSSIA API    
    \item Many "low-hanging fruit" TODOs
    \item Mobile devices ports : 
    \begin{itemize}
        \item Android : builds and run but requires adapted UI.
        \item Web port : with PNaCl, runs but crashes. Will open the way to WebAssembly. 
        \item iDevices (many artists use them).
    \end{itemize}
\end{itemize}
    
\end{frame}



\begin{frame}
    \frametitle{Links}
    \begin{itemize}
        \item \textbf{Grab a release !} :~\\ \url{github.com/OSSIA/i-score/releases}. 
        \item \textbf{Protocols and implementations} :~\\
        \url{github.com/OSSIA}
    \end{itemize}
        
    \centering
    \vspace{2cm}
    \Large{Thanks ! Questions ?}
    \vspace{2cm}
    
    \small{Credits: 'simple' Beamer theme, Facundo Muñoz; Fira font}
\end{frame}    
\end{document}
